% Progress Report Template (LaTeX)
% 可直接编译：pdflatex report.tex
% 将示例文本替换成你的内容即可使用

\documentclass[11pt,a4paper]{article}
\usepackage[margin=2.5cm]{geometry}
\usepackage{graphicx}
\usepackage{hyperref}
\usepackage{enumitem}

% 中文用户取消下面两行注释
% \usepackage{ctex}
% \setCJKmainfont{Noto Sans CJK SC}

\title{Weekly Progress Report \\ \large 2026-03-25 to 2026-04-01}
\author{Your Name}
\date{}

\begin{document}
\maketitle

\section{Summary}
This week I validated the mixup direction on CIFAR-10, collected evidence from code and result artifacts, and narrowed down the next experiment choices.

\section{Progress}

\subsection{Experiment A}
Tested mixup on CIFAR-10 after last week's meeting. Accuracy improved from 78.3\% to 81.7\%. Evidence includes commit \texttt{abc1234} and the result figure in \texttt{results/example.png}.

% 插入结果图（取消注释并替换路径）：
% \begin{figure}[h]
%   \centering
%   \includegraphics[width=0.8\textwidth]{results/example.png}
%   \caption{实验结果}
% \end{figure}

\subsection{Reading and Discussion}
Read the CutMix paper, summarized three follow-up ideas, and discussed the detection-task extension with a senior labmate.

\section{Questions for Advisor}
\begin{itemize}[leftmargin=2em]
  \item Which follow-up direction should go first: adaptive mixup, contrastive learning, or detection?
  \item Should we request extra GPU time for the ImageNet run?
\end{itemize}

\section{Blockers}
The GPU queue is long, so the ImageNet experiment may be delayed by about a week.

\section{Next Steps}
\begin{itemize}[leftmargin=2em]
  \item Validate the result on an ImageNet subset.
  \item Start drafting the related-work section.
\end{itemize}

\end{document}
